\documentclass[a4paper,10pt]{report}\input{../0.Préambule/Préambule_cours_prof}\begin{document}

\newpage
\section*{La symétrie axiale dans la vie de tous les jours}
\addcontentsline{toc}{section}{La symétrie axiale dans la vie de tous les jours}

\subsection*{Le jeu des 7 différences}
\addcontentsline{toc}{subsection}{Le jeu des 7 différences}

La figure ci-contre devrait posséder un axe de symétrie mais des erreurs se sont glissées.

\medskip
\noindent Combien en trouves-tu ?

\begin{center}

\includegraphics[scale=0.55]{../40.Image/13. Act-La symétrie axiale 1.png}
\end{center}

\subsection*{En attendant l'ASSR!}
\addcontentsline{toc}{subsection}{En attendant l'ASSR!}

Sur les panneaux suivants, trace le ou les axes de symétrie, s'il(s) existe(nt) :

\begin{center}
\includegraphics[scale=0.9]{../40.Image/13. Act-La symétrie axiale 2.png}
\end{center}
\end{document}